{
\def\AMCbeginQuestion#1#2{\par\noindent{\bf (POSCOMP 2004 - Questão 68)}}
\begin{question}{} O protocolo padrão para gerenciamento de redes TCP/IP, definido pelo IETF, é:
\begin{choices}
\correctchoice{SNMP}
\wrongchoice{SMTP}
\wrongchoice{HTTP}
\wrongchoice{COPS}
\wrongchoice{SSH}
\end{choices}
\end{question}
}
%
{
\def\AMCbeginQuestion#1#2{\par\noindent{\bf (POSCOMP 2004 - Questão 69)}}
\begin{question}{} Qual das opções abaixo melhor caracteriza o protocolo IP?
\begin{choices}
\correctchoice{Não orientado a conexão, sem suporte a QoS, sem mecanismo de retransmissão.}
\wrongchoice{Orientado a conexão, com suporte a QoS, com mecanismo de retransmissão.}
\wrongchoice{Orientado a conexão, sem suporte a QoS, sem mecanismo de retransmissão.}
\wrongchoice{Orientado a conexão, sem suporte a QoS, com mecanismo de retransmissão.}
\wrongchoice{Não orientado a conexão, com suporte a QoS, sem mecanismo de retransmissão.}
\end{choices}
\end{question}
}
%
{
\def\AMCbeginQuestion#1#2{\par\noindent{\bf (POSCOMP 2004 - Questão 70)}}
\begin{question}{} Assinale a alternativa que apresenta um protocolo de roteamento baseado no algoritmo vetor-distância e é classificado como IGP (\textit{Interior Gateway Protocol}):
\begin{choices}
\correctchoice{RIP}
\wrongchoice{OSPF}
\wrongchoice{ICMP}
\wrongchoice{BGP}
\wrongchoice{RSVP}
\end{choices}
\end{question}
}
%
{
\def\AMCbeginQuestion#1#2{\par\noindent{\bf (POSCOMP 2005 - Questão 63)}}
\begin{question}{} Os protocolos de transporte atribuem a cada serviço um identificador único, o qual é empregado para encaminhar uma requisição de um aplicativo cliente ao processo servidor correto. Nos protocolos de transporte TCP e UDP, como esse identificador se denomina?
\begin{choices}
\correctchoice{Porta.}
\wrongchoice{Endereço IP.}
\wrongchoice{Conexão.}
\wrongchoice{Identificador do processo (PID).}
\wrongchoice{Protocolo de aplicação.}
\end{choices}
\end{question}
}
%
{
\def\AMCbeginQuestion#1#2{\par\noindent{\bf (POSCOMP 2005 - Questão 64)}}
\begin{question}{} O DNS (\textit{Domain Name System}) é um serviço de diretórios na Internet que:
\begin{choices}
\correctchoice{Traduz o nome de um hospedeiro (host) para seu endereço IP.}
\wrongchoice{Localiza a instituição à qual um dado host pertence.}
\wrongchoice{Retorna a porta da conexão TCP do host.}
\wrongchoice{Retorna a porta da conexão UDP do host.}
\wrongchoice{Traduz o endereço IP de um hospedeiro para um nome de domínio na Internet.}
\end{choices}
\end{question}
}
%
{
\def\AMCbeginQuestion#1#2{\par\noindent{\bf (POSCOMP 2005 - Questão 65)}}
\begin{question}{} Um dos mecanismos de congestionamento na rede é o que utiliza temporizadores de transmissão e duas variáveis chamadas de: Janela de Congestionamento e Patamar. A Janela de Congestionamento impõe um limite à quantidade de tráfego que um host pode enviar dentro de uma conexão. O Patamar é uma variável que regula o crescimento da Janela de Congestionamento durante as transmissões daquela conexão. Assinale a alternativa correta:
\begin{choices}
\correctchoice{A Janela de Congestionamento dobra de tamanho (cresce exponencialmente) quando a confirmação das mensagens enviadas ocorre antes dos temporizadores de retransmissão se esgotarem (\textit{time-out}), até o limite do Patamar.}
\wrongchoice{A quantidade de mensagens não confirmadas na transmissão, num dado instante, deve ser superior ao mínimo entre a Janela de Congestionamento e a Janela de Recepção desta conexão.}
\wrongchoice{Após exceder o valor de Patamar ainda sem esgotar os temporizadores, a janela decresce linearmente.}
\wrongchoice{Quando excede o valor de Patamar e esgotam os temporizadores, a janela decresce exponencialmente.}
\wrongchoice{Todas as alternativas estão corretas.}
\end{choices}
\end{question}
}
%
{
\def\AMCbeginQuestion#1#2{\par\noindent{\bf (POSCOMP 2005 - Questão 66)}}
\begin{question}{} Algoritmos de roteamento são o meio que um roteador utiliza para encaminhar mensagens na camada de rede. Assinale a alternativa incorreta.
\begin{choices}
\correctchoice{No roteamento de Estado de Enlace (\textit{Link State}), os valores dos enlaces são calculados pelo projetista da rede e os roteadores atualizam suas tabelas por estes valores.}
\wrongchoice{Nos algoritmos de roteamento estáticos as rotas são determinadas via tabelas definidas a priori e fixadas para o roteador, em geral manualmente.}
\wrongchoice{No roteamento por Vetor de Distância (\textit{Distance Vector}), as tabelas de roteamento definidas pelos roteadores vizinhos são repassadas periodicamente a cada roteador para obtenção de sua própria tabela.}
\wrongchoice{Algoritmos de roteamento buscam estabelecer o caminho de menor custo entre dois hosts através do cálculo dos custos acumulados mínimos entre os enlaces disponíveis, dada a topologia da rede.}
\wrongchoice{O OSPF é um exemplo de protocolo de roteamento baseado em Estado de Enlace e o BGP é um exemplo de protocolo de roteamento baseado em Vetor de Distâncias.}
\end{choices}
\end{question}
}
%
{
\def\AMCbeginQuestion#1#2{\par\noindent{\bf (POSCOMP 2005 - Questão 67)}}
\begin{question}{} Sejam as afirmações:
\begin{enumerate}[label=\Roman*.]
\item O HTTP e o FTP são protocolos da camada de aplicação e utilizam o protocolo de transporte TCP.
\item Ambos (HTTP e FTP) utilizam duas conexões TCP, uma para controle da transferência e outra para envio dos dados transferidos (controle fora da banda).
\item O HTTP pode usar conexões não persistentes e persistentes. O HTTP/1.0 usa conexões não persistentes. O modo default do HTTP/1.1 usa conexões persistentes.
\end{enumerate}
Dadas estas três afirmações, indique qual a alternativa correta:
\begin{choices}
\correctchoice{Somente I e III são verdadeiras.}
\wrongchoice{I, II e III são verdadeiras.}
\wrongchoice{Somente I e II são verdadeiras.}
\wrongchoice{Somente I é verdadeira.}
\wrongchoice{I, II e III são falsas.}
\end{choices}
\end{question}
}
%
{
\def\AMCbeginQuestion#1#2{\par\noindent{\bf (POSCOMP 2006 - Questão 64)}}
\begin{question}{} Sobre o protocolo IP (\textit{Internet Protocol}), é correto afirmar:
\begin{choices}
\correctchoice{O espaço de endereçamento do IPv4 e do IPv6 é de 32 e 128 bits, respectivamente.}
\wrongchoice{O tamanho do cabeçalho do IPv4 é fixado em 96 bits.}
\wrongchoice{O cabeçalho IP inclui informação sobre o protocolo de camada de enlace empregado.}
\wrongchoice{A classe C de endereços IPv4 reserva 16 bits para endereço de rede.}
\wrongchoice{O roteamento IP associa o endereço IP com o número de porta em nível de transporte.}
\end{choices}
\end{question}
}
%
{
\def\AMCbeginQuestion#1#2{\par\noindent{\bf (POSCOMP 2006 - Questão 65)}}
\begin{question}{} Duas tecnologias utilizadas para acesso residencial à Internet são ADSL e Cable Modem. Qual afirmação é incorreta?
\begin{choices}
\correctchoice{Os canais de \textit{upstream} e \textit{downstream} da tecnologia Cable Modem necessitam de contenção de acesso.}
\wrongchoice{Ambas permitem taxas de transmissão diferentes para \textit{upstream} e \textit{downstream}.}
\wrongchoice{Os canais de \textit{upstream} e \textit{downstream} da tecnologia ADSL não necessitam de contenção de acesso.}
\wrongchoice{ADSL utiliza par trançado dedicado para cada residência.}
\wrongchoice{Cable Modem utiliza cabo compartilhado para diversas residências.}
\end{choices}
\end{question}
}
%
{
\def\AMCbeginQuestion#1#2{\par\noindent{\bf (POSCOMP 2006 - Questão 66)}}
\begin{question}{} Os endereços IP são divididos em classes. Qual afirmação é incorreta?
\begin{choices}
\correctchoice{Máscaras podem dividir o campo Rede do endereço IP em Rede e Sub-rede para facilitar o roteamento interno.}
\wrongchoice{Existem mais redes classe B do que classe A.}
\wrongchoice{Uma rede classe C permite mais hosts do que uma rede classe B.}
\wrongchoice{A classe D é dedicada a endereços \textit{multicast}.}
\wrongchoice{NAT (Tradução de Endereço de Rede) é utilizada em redes com vários hosts que se conectam à Internet através de poucos endereços IP.}
\end{choices}
\end{question}
}
%
{
\def\AMCbeginQuestion#1#2{\par\noindent{\bf (POSCOMP 2007 - Questão 64)}}
\begin{question}{} O controle de congestionamento é uma das funções desempenhadas pela Camada de Transporte no modelo TCP/IP. Sobre essa função, assinale a alternativa INCORRETA.
\begin{choices}
\correctchoice{No controle de congestionamento assistido pela rede, os nodos (roteadores) enviam notificações explícitas do estado de congestionamento da rede diretamente à fonte de cada fluxo que, por meio dele, trafega.}
\wrongchoice{No controle de congestionamento fim-a-fim, uma situação de congestionamento é intuída pelos hosts terminais via eventos como perda ou atraso excessivo de pacotes.}
\wrongchoice{O mecanismo \textit{Explicit Congestion Notification} (ECN) utiliza um dos dois últimos bits do campo ToS do cabeçalho IPv4 para notificar a um destinatário o estado de congestionamento da rede.}
\wrongchoice{Ao perceber um estado de congestionamento na rede, uma conexão TCP, por meio de seu mecanismo de prevenção de congestionamento (\textit{congestion avoidance}), reduz o tamanho de sua janela de congestionamento.}
\wrongchoice{Na fase de partida lenta (\textit{slow start}) de uma conexão TCP, o tamanho da janela de congestionamento aumenta a cada RTT (\textit{Round-Trip Time}) de forma exponencial, até que esse tamanho alcance um determinado valor de limiar (\textit{threshold}).}
\end{choices}
\end{question}
}
%
{
\def\AMCbeginQuestion#1#2{\par\noindent{\bf (POSCOMP 2007 - Questão 65)}}
\begin{question}{} Sobre o protocolo de transferência de hipertextos (HTTP - \textit{Hyper-Text Transfer Protocol}), é CORRETO afirmar que
\begin{choices}
\correctchoice{A instrução GET condicional permite que o cliente opte por receber um determinado objeto do servidor apenas se este tiver sido alterado depois de uma determinada data e hora.}
\wrongchoice{O protocolo HTTP é capaz de transportar nativamente arquivos no formato binário.}
\wrongchoice{A versão 1.0 do protocolo HTTP não permite a utilização de cookies.}
\wrongchoice{A versão 1.1 do protocolo HTTP difere da versão 1.0 na capacidade de transportar objetos maiores.}
\wrongchoice{O protocolo HTTP não pode ser utilizado para transportar outros tipos de objetos senão os hiper-textos.}
\end{choices}
\end{question}
}
%
{
\def\AMCbeginQuestion#1#2{\par\noindent{\bf (POSCOMP 2007 - Questão 66)}}
\begin{question}{} Considere os pares de endereços de hosts e suas respectivas máscaras de endereços listados abaixo.
\begin{enumerate}[label=\Roman*.]
\item 192.168.0.43/255.255.255.192 e 192.168.0.66/255.255.255.192
\item 192.168.1.97/255.255.255.224 e 192.168.1.118/255.255.255.224
\item 192.168.2.115/255.255.255.128 e 192.168.2.135/255.255.255.128
\item 192.168.3.34/255.255.255.240 e 192.168.3.46/255.255.255.240
\item 192.168.4.167/255.255.255.224 e 192.168.4.207/255.255.255.224
\end{enumerate}
Os itens nos quais o par citado pertence a uma mesma sub-rede são
\begin{choices}
\correctchoice{apenas II, IV.}
\wrongchoice{apenas I, II, V.}
\wrongchoice{apenas I, III.}
\wrongchoice{apenas II, III, IV.}
\wrongchoice{apenas III, IV, V.}
\end{choices}
\end{question}
}
%
{
\def\AMCbeginQuestion#1#2{\par\noindent{\bf (POSCOMP 2007 - Questão 67)}}
\begin{question}{} Analise as seguintes afirmativas.
\begin{enumerate}[label=\Roman*.]
\item O protocolo UDP é um protocolo da Camada de Transporte orientado a datagrama, enquanto que o TCP é um protocolo da Camada de Transporte orientado a conexão.
\item Apesar de o protocolo IP ser orientado a datagrama, o protocolo UDP é necessário por fornecer multiplexação de um endereço de rede em várias portas, permitindo que múltiplos processos sejam endereçados em um mesmo endereço de rede.
\item O protocolo TCP utiliza o tamanho da janela deslizante de uma conexão para o controle de congestionamento.
\end{enumerate}
A esse respeito, pode-se afirmar que
\begin{choices}
\correctchoice{todas as afirmativas são corretas.}
\wrongchoice{somente a afirmativa I é correta.}
\wrongchoice{somente as afirmativas I e II são corretas.}
\wrongchoice{somente as afirmativas I e III são corretas.}
\wrongchoice{somente as afirmativas II e III são corretas.}
\end{choices}
\end{question}
}
%
{
\def\AMCbeginQuestion#1#2{\par\noindent{\bf (POSCOMP 2008 - Questão 07)}}
\begin{question}{} Analise as seguintes afirmativas.
\begin{enumerate}[label=\Roman*.]
\item Um servidor DNS suporta dois tipos de consulta: iterativa e recursiva. Na consulta iterativa que é a mais utilizada, caso um servidor DNS não tenha a informação pedida pela máquina solicitante, ele irá buscar a mesma consultando outros servidores.
\item Como estratégia para aumentar a confiabilidade na resposta dos servidores DNS quando do emprego de caches, devem-se utilizar valores grandes de TLL (\textit{Time-To-Live}), mantendo elevado o tempo de validade do registro na cache.
\item Um servidor DNS pode atender dois tipos de consultas: tradução direta, na qual, a partir de um endereço IP, o mesmo retorna o nome de rede do equipamento; e tradução inversa, na qual, a partir de um nome de rede, retorna o IP associado ao mesmo.
\end{enumerate}
A análise permite concluir que
\begin{choices}
\correctchoice{nenhuma afirmativa está correta.}
\wrongchoice{somente a afirmativa I está correta.}
\wrongchoice{somente a afirmativa II está correta.}
\wrongchoice{somente a afirmativa III está correta.}
\wrongchoice{todas as afirmativas estão corretas.}
\end{choices}
\end{question}
}
%
{
\def\AMCbeginQuestion#1#2{\par\noindent{\bf (POSCOMP 2008 - Questão 08)}}
\begin{question}{} O nível de transporte oferece serviços para comunicação entre computadores, independentemente das tecnologias utilizadas nos outros níveis.

Analise as seguintes afirmativas relativas à confirmação do recebimento de pacotes no nível de transporte.
\begin{enumerate}[label=\Roman*.]
\item A troca de dados entre um computador transmissor e um receptor não precisa obrigatoriamente de uma confirmação para cada pacote enviado. Existem três estratégias que podem ser utilizadas: confirmação seletiva, confirmação cumulativa e confirmação em bloco.
\item Na confirmação seletiva, cada pacote recebido por um computador não gera uma informação de confirmação individualizada para o computador que enviou o pacote.
\item Na confirmação do recebimento de pacotes, o consumo da banda de rede pode ser otimizado pelo uso de um mecanismo denominado \textit{piggybacking}. No \textit{piggybacking} a informação de confirmação ``pega carona'' em mensagem de dados que retorna ao computador emissor como consequência do fluxo normal de troca de dados.
\end{enumerate}
A análise permite concluir que
\begin{choices}
\correctchoice{apenas as afirmativas I e III estão corretas.}
\wrongchoice{nenhuma das afirmativas está correta.}
\wrongchoice{apenas as afirmativas I e II estão corretas.}
\wrongchoice{apenas as afirmativas II e III estão corretas.}
\wrongchoice{todas as afirmativas estão corretas.}
\end{choices}
\end{question}
}
%
{
\def\AMCbeginQuestion#1#2{\par\noindent{\bf (POSCOMP 2008 - Questão 14)}}
\begin{question}{} Sobre a comunicação entre processos distribuídos, é CORRETO afirmar
\begin{choices}
\correctchoice{que, nos sockets do tipo datagrama, o endereço do socket do processo correspondente acompanha cada envio de mensagem como um parâmetro da primitiva \texttt{sendto()}.}
\wrongchoice{que, no modo síncrono de envio de mensagem, o processo que recebe a mensagem terá sua execução desviada por uma interrupção de sistema operacional para tratar uma mensagem recebida.}
\wrongchoice{que um processo tem no máximo uma porta (port) para receber as mensagens dos seus interlocutores.}
\wrongchoice{que multiportas são estruturas do tipo portas multicast que permitem a comunicação N x M entre processos.}
\wrongchoice{que, nos sockets do tipo datagrama, a primitiva \texttt{sNew()} é usada para aceitar a conexão solicitada por um processo que solicita comunicação.}
\end{choices}
\end{question}
}
%
